\documentclass[12pt]{article}
\usepackage[utf8]{inputenc}
\usepackage[T1]{fontenc}
\usepackage{hyperref}
\title{Purig Language Resources for K-12 Curriculum Development}
\author{Compiled by Claude Sonnet 5.5 (claude-sonnet-5-5) for Andr\'as Kornai}
\date{}
\begin{document}
\maketitle

\section{Introduction}
Purig (Purgi, Purik; ISO 639-3 \texttt{prx}) is a Tibetic language closely related to Ladakhi--Balti, spoken mainly in the Kargil area of Ladakh and in Baltistan. Wikipedia cites about 94,000 native speakers in India (2011 census), higher than the 38,000 in the survey prompt. Linguists disagree on whether it is distinct from Balti. Advocates in Ladakh have called for Purigi to be a medium of instruction in schools.

\section{Primary Educational Resources}
No textbooks or primers were found. A policy item concerning the medium of instruction:
\begin{itemize}
\item AJUIAK demand for Purigi as a medium of instruction in schools: \url{https://voiceofladakh.in/?p=1690}
\end{itemize}

\section{Linguistic and Reference Materials}
\begin{itemize}
\item Wikipedia, Purgi language: \url{https://en.wikipedia.org/wiki/Purgi_language}
\item Purgi Dictionary app (iPad/iOS), by Dr.\ Winsy Abraham, over 5000 words, English to Purgi with Hindi and Urdu, Latin script, authors warn of errors: \url{https://apps.apple.com/app/id6780769907}
\item Zemp, M., \emph{A Grammar of Purik Tibetan} (Brill, 2018), cited in the Wikipedia references. Its scope was not verified, and no direct URL was seen.
\end{itemize}

\section{Religious and Community Resources}
\begin{itemize}
\item Voice of Ladakh, ``What do Experts tell About Purigi Language?'': \url{https://voiceofladakh.in/2021/03/what-do-experts-tell-about-purigi-language/}
\end{itemize}

\section{Digital and Audio Resources}
The Purgi Dictionary app above is the only digital resource found. No audio or video was found.

\section{Curriculum Development Context}
There is advocacy for school use, but no confirmed curriculum materials were found. A further search for textbooks and the Purig script, and contact with Kargil education bodies, would be needed.

\end{document}